%%%%%%%%%%%%%%%%%%%%%%%%%%%%%%%%%%%%%%%%%%%%%%%%%%%%%%%%%%%%%%%%%%%%%%
% writeLaTeX Example: A quick guide to LaTeX
%
% Source: Dave Richeson (divisbyzero.com), Dickinson College
% 
% A one-size-fits-all LaTeX cheat sheet. Kept to two pages, so it 
% can be printed (double-sided) on one piece of paper
% 
% Feel free to distribute this example, but please keep the referral
% to divisbyzero.com
% 
%%%%%%%%%%%%%%%%%%%%%%%%%%%%%%%%%%%%%%%%%%%%%%%%%%%%%%%%%%%%%%%%%%%%%%
% How to use writeLaTeX: 
%
% You edit the source code here on the left, and the preview on the
% right shows you the result within a few seconds.
%
% Bookmark this page and share the URL with your co-authors. They can
% edit at the same time!
%
% You can upload figures, bibliographies, custom classes and
% styles using the files menu.
%
% If you're new to LaTeX, the wikibook is a great place to start:
% http://en.wikibooks.org/wiki/LaTeX
%
%%%%%%%%%%%%%%%%%%%%%%%%%%%%%%%%%%%%%%%%%%%%%%%%%%%%%%%%%%%%%%%%%%%%%%

\documentclass[10pt,landscape]{article}
\usepackage{amssymb,amsmath,amsthm,amsfonts}
\usepackage{multicol,multirow}
\usepackage{calc}
\usepackage{ifthen}
\usepackage{siunitx}
\usepackage{amsmath}
\usepackage{amsthm}
\usepackage{mathtools}
\usepackage{graphicx}
\usepackage{enumitem}
\usepackage{nicematrix}
\usepackage{color} 
\usepackage{cancel}
\usepackage{longdivision}
\usepackage{xfrac}
\usepackage[landscape]{geometry}
\usepackage[colorlinks=true,citecolor=blue,linkcolor=blue]{hyperref}

% Contains all the "\newcommand" commands
% Colored Text
\newcommand{\redtext}[1]{\textcolor{red}{#1}}
\newcommand{\orangetext}[1]{\textcolor{orange}{#1}}
\newcommand{\yellowtext}[1]{\textcolor{yellow}{#1}}
\newcommand{\greentext}[1]{\textcolor{green}{#1}}
\newcommand{\bluetext}[1]{\textcolor{blue}{#1}}
\newcommand{\purpletext}[1]{\textcolor{purple}{#1}}
\newcommand{\pinktext}[1]{\textcolor{pink}{#1}}
\newcommand{\violettext}[1]{\textcolor{violet}{#1}}

% Highlighted Text
\newcommand{\redbox}[1]{\colorbox{red}{#1}}
\newcommand{\orangebox}[1]{\colorbox{orange}{#1}}
\newcommand{\yellowbox}[1]{\colorbox{yellow}{#1}}
\newcommand{\greenbox}[1]{\colorbox{green}{#1}}
\newcommand{\bluebox}[1]{\colorbox{blue}{\textcolor{white}{#1}}}
\newcommand{\purplebox}[1]{\colorbox{purple}{\textcolor{white}{#1}}}
\newcommand{\pinkbox}[1]{\colorbox{pink}{#1}}
\newcommand{\violetbox}[1]{\colorbox{violet}{#1}}

%%% Center Text %%%
\newcommand{\centertext}[1]
{%
    \begin{center}
        #1
    \end{center}
}

% ===== Tables =====

% Make a table with a header
% #1 is |c|c|c|
% #2 is H1 & H2 & H3 \\
% Body is B1 & B2 & B3 \\
\newenvironment{tablewithheader}[2]
{% begin
    \begin{center}

        \begin{tabular}{#1}

            % Header
            \hline
            #2 \\ [0.5ex]
            \hline

            % Body
            % =====
            \hline

}
{% end

            \hline
            % =====

        \end{tabular}
    \end{center}
}

% Make a table with a header aligned to the left
% #1 is |c|c|c|
% #2 is H1 & H2 & H3 \\
% Body is B1 & B2 & B3 \\

% EXAMPLE
% \begin{tablewithheader}
% { |c|c|c| }
% { Col A & Col B & Col C }

%     1 & 2 & 3 \\
%     a & b & c \\ 
%     x & y & z \\

% \end{tablewithheader}

\newenvironment{tablewithheaderleft}[2]
{% begin

    \begin{tabular}{#1}

        % Header
        \hline
        #2 \\ [0.5ex]
        \hline

        % Body
        % =====
        \hline

}
{% end

        \hline
        % =====

    \end{tabular}
}

% Make a table without a header
% #1 is |c|c|c|
% Body is B1 & B2 & B3 \\
\newenvironment{tablewithoutheader}[1]
{% begin
    \begin{center}

        \begin{tabular}{#1}

            % Body
            % =====
            \hline

}
{% end

            \hline
            % =====

        \end{tabular}
    \end{center}
}

% FFFF00 (yellow)
% FFFF9E (yellow)
% When using |c|c|c| replace "c" with "s" for a highlighted column
\newcolumntype{s}{>{\columncolor[HTML]{FFFF9E}} c}

% ===== Lists =====

% Numbered List
% 
% \begin{enumerate}
%     \item ...
% \end{enumerate}

% Bullet-point List
% 
% \begin{itemize}
%     \item ...
% \end{itemize}

% Alphabet List
\newlist{alphabet}{enumerate}{1}
\setlist[alphabet]{label=(\alph*)}

% Roman Numeral List
\newlist{romanlist}{enumerate}{1}
\setlist[romanlist]{label=(\roman*)}

% ===== Math =====

% Multi-lined Math
% 
%   - Everything aligns at &
% 
% \begin{align*}
%     a + b &= c \\
%     &= 2 + 2 \tag{ This is a tag! } \\
%     &= 4 \\
% \end{align*}

% ========== Set Notation ==========

% Italics Shorthand
\newcommand{\mathset}[1]
{
    \ensuremath{\mathit{#1}}
}

% Brackets Shorthand
\newcommand{\mathinset}[1]
{
    \ensuremath{\{ #1 \}}
}

% Define colors for set notation
\newcommand{\definesetcolors}{%

    % color of set of Natural numbers
    \colorlet{colN}{Red!60}
    % color of set of Whole Numbers
    \colorlet{colW}{OrangeRed!60}
    % color of set of Integers
    \colorlet{colZ}{Orange!60}
    % color of set of Rational numbers
    \colorlet{colQ}{Gold!60}
    % color of set of Real numbers
    \colorlet{colR}{ForestGreen!50}
    % color of set of Complex numbers
    \colorlet{colC}{CornflowerBlue!80}
    % color of set of Hypercomplex numbers
    \colorlet{colH}{BlueViolet!50}
    % color of set of Transcendental numbers
    % \colorlet{colT}{Green!30}

    % center coordinate of Natural numbers
    \coordinate (oN) at (0.0, 0.0);
    % center coordinate of Whole Numbers
    \coordinate (oW) at (0.5, 0.5);
    % center coordinate of Integers
    \coordinate (oZ) at (1.0, 1.0);
    % center coordinate of Rational numbers
    \coordinate (oQ) at (1.5, 1.5);
    % center coordinate of Real numbers
    \coordinate (oR) at (2.0, 2.0);
    % center coordinate of Complex numbers
    \coordinate (oC) at (2.5, 2.5);

    % set of Natural numbers
    \def\setN{(oN) ellipse (2.5 and 1.0)}
    % set of Whole numbers
    \def\setW{(oW) ellipse (3.5 and 2.0)}
    % set of Integers numbers
    \def\setZ{(oZ) ellipse (4.5 and 3.0)}
    % set of Rational numbers
    \def\setQ{(oQ) ellipse (5.5 and 4.0)}
    % set of Real numbers
    \def\setR{(oR) ellipse (6.5 and 5.0)}
    % set of Complex numbers
    \def\setC{(oC) ellipse (7.5 and 6.0)}
}

% ========== Discrete Specific ==========

\newcommand{\T}{\greentext{T}}
\newcommand{\F}{\redtext{F}}

\newcommand{\TRUE}{\greentext{TRUE}}
\newcommand{\FALSE}{\redtext{FALSE}}

% Contains misc fonts

% Common Sets
\newcommand{\setN}{\ensuremath{\mathbb{N}}}
\newcommand{\setZ}{\ensuremath{\mathbb{Z}}}
\newcommand{\setQ}{\ensuremath{\mathbb{Q}}}
\newcommand{\setR}{\ensuremath{\mathbb{R}}}
\newcommand{\setC}{\ensuremath{\mathbb{C}}}





\ifthenelse{\lengthtest { \paperwidth = 11in}}
    { \geometry{top=.5in,left=.5in,right=.5in,bottom=.5in} }
	{\ifthenelse{ \lengthtest{ \paperwidth = 297mm}}
		{\geometry{top=1cm,left=1cm,right=1cm,bottom=1cm} }
		{\geometry{top=1cm,left=1cm,right=1cm,bottom=1cm} }
	}
\pagestyle{empty}
\makeatletter
\renewcommand{\section}{\@startsection{section}{1}{0mm}%
                                {-1ex plus -.5ex minus -.2ex}%
                                {0.5ex plus .2ex}%x
                                {\normalfont\large\bfseries}}
\renewcommand{\subsection}{\@startsection{subsection}{2}{0mm}%
                                {-1explus -.5ex minus -.2ex}%
                                {0.5ex plus .2ex}%
                                {\normalfont\normalsize\bfseries}}
\renewcommand{\subsubsection}{\@startsection{subsubsection}{3}{0mm}%
                                {-1ex plus -.5ex minus -.2ex}%
                                {1ex plus .2ex}%
                                {\normalfont\small\bfseries}}
\makeatother
\setcounter{secnumdepth}{0}
\setlength{\parindent}{0pt}
\setlength{\parskip}{0pt plus 0.5ex}
% -----------------------------------------------------------------------

\begin{document}

\raggedright
\footnotesize

\begin{multicols}{4}
\setlength{\premulticols}{1pt}
\setlength{\postmulticols}{1pt}
\setlength{\multicolsep}{1pt}
\setlength{\columnsep}{2pt}

% -----------------------------------------------------------------------
% PERMUTATIONS
% -----------------------------------------------------------------------
\section{Permutations}

    The number of ways to arrange $r$ objects out of $n$ is where \textbf{order matters} is...

    The number of $r$ length words over an alphabet with $n$ letters \textbf{without} repetitions

    \begin{align*}
        n: \; & \text{Total \# of objects} \\
        r: \; & \text{\# of objects you want to choose from} \\
        P(n, r) = \prescript{}{n}{P}_r &= n(n - 1)(n - 2) \cdot \ldots \cdot (n - r + 1) \\
        &= \frac{n!}{(n - r)!} \\
    \end{align*}

    Some examples are:

    \begin{itemize}
        \item Different ways of choosing r roles out of a class of n students, where each role has one student playing it, and no students have multiple roles
        \item Different ways of passing out 5 cards to a group of 5 people out of a deck of 52, where each person gets one card
        \item Different ways of having 1st 2nd 3rd place runners in a race of 10 athletes.
        \item Different ways to make a 4-letter word using different letters.
    \end{itemize}

% -----------------------------------------------------------------------
% COMBINATIONS
% -----------------------------------------------------------------------
\section{Combinations}

\begin{enumerate}
    \item How many ways can I pick $r$ objects from $n$ distinct objects in an \textbf{unordered} set?
    \item $r$-element subsets of a set of cardinality $n$
    \item Length $n$ binary strings of length $r$
\end{enumerate}

\begin{align*}
    C(n, r) &= \prescript{}{n}{C}_r = {n \choose r} = \frac{n!}{r! \cdot (n - r)!}
\end{align*}

% -----------------------------------------------------------------------
% STARS AND BARS
% -----------------------------------------------------------------------
\section{Stars and Bars}

\begin{align*}
    & {n + k - 1 \choose k - 1} = {n + k - 1 \choose n} \\
    &n: \text{ number of items} \\
    &k: \text{ number of bins} \\
\end{align*}

Where 

\begin{enumerate}
    \item Order doesn't matter
    \item Bins can contain 0
    \item All items must be placed
\end{enumerate}

% -----------------------------------------------------------------------
% INTEGER EQUATIONS
% -----------------------------------------------------------------------
\section{Integer Equations}

\begin{align*}
    & {n + k - 1 \choose k - 1} = {n + k - 1 \choose n} \\
    &n: \text{ RHS } \\
    &k: \text{ Number of variables } \\
    &\\
    & {n + k(1 - m) - 1 \choose k - 1} \\
    &n: \text{ Number of items } \\
    &k: \text{ Number of bins } \\
    &m: \text{ At Least } \\
\end{align*}

For 

$$(a_1 + c_1) + (a_2 + c_2) + \ldots (a_k + c_k) = n$$

This is just saying that $a_i \geq c_i$, so just move all the constants to the left and solve like normal

$$a_1 + a_2 + \ldots + a_k = n - c_1 - c_2 - \ldots - c_k$$

% -----------------------------------------------------------------------
% BINOMIAL THEOREM
% -----------------------------------------------------------------------
\section{Binomial Theorem}
$$\text{For } n = 0, 1, 2, \ldots, (a + b)^n = \sum_{k = 0}^{n} {n \choose k} a^{n - k}b^k$$

% -----------------------------------------------------------------------
% INCLUSION EXCLUSION (PIE)
% -----------------------------------------------------------------------
\section{Inclusion Exclusion (PIE)}
Let $A_1, A_2, \ldots , A_n$ be finite sets. Then

\begin{align*}
    &|A_1 \cup A_2 \cup \ldots \cup A_n| = \sum_{1 \leq i \leq n}^{} |A_i| \\
    &- \sum_{1 \leq i < j \leq n}^{} | A_i \cap A_j | \\
    &+ \sum_{1 \leq i < j < k \leq n}^{} |A_i \cap A_j \cap A_k | \\
    &- \ldots + (-1)^{n + 1} |A_1 \cap A_2 \cap \ldots \cap A_n | \\
\end{align*}

\begin{center}
    \textbf{ODD} sets involved: \textbf{ADD} \\
    \textbf{EVEN} sets involved: \textbf{SUBTRACT} \\
\end{center}

% -----------------------------------------------------------------------
% CARDINALITY OF INTERSECTIONS
% -----------------------------------------------------------------------
\section{Cardinality of Intersections}
Let $A_1, A_2, \ldots , A_n$ be finite sets. Then

\begin{align*}
    &|A_1 \cap A_2 \cap \ldots \cap A_n | = |U| - \sum_{i = 1}^{n} |\overline{A_i}| \\
    &+ \sum_{k_1 < k_2}^{} |\overline{A_{k_1}} \cap \overline{A_{k_2}}| \\
    &- \sum_{k_1 < k_2 < k_3}^{} |\overline{A_{k_1}} \cap \overline{A_{k_2}} \cap \overline{A_{k_3}}| \\
    &+ \ldots + (-1)^n |\overline{A_{k_1}} \cap \overline{A_{k_2}} \cap \ldots \cap \overline{A_n}|
\end{align*}

The big idea is

\begin{center}
    \textbf{ODD} number of sets, then \textbf{SUBTRACT} \\
    \textbf{EVEN} number of sets, then \textbf{ADD} \\
\end{center}

% -----------------------------------------------------------------------
% PROBLEMS
% -----------------------------------------------------------------------
\section{Problems}

% -----------------------------------------------------------------------
% LICENSE PLATES
% -----------------------------------------------------------------------
\subsection{License Plates}

How many 7-character license plates are there that are all different letters and they are
in alphabetical order?

$$26 \choose 7$$

How many 7-character license plates are there from all capital letters and they are
in alphabetical order? They can be duplicates, so AAAAAAA is valid

Use stars and bars! $x_A, x_B, x_C$ represents the number  of appearances of each letter. They all
must sum to 7 

$n = 7,\; k = 26$

$${7 + 26 - 1 \choose 7} = {32 \choose 7}$$

In general

\begin{align*}
    n &= \text{string length} \\
    k &= \text{size of alphabet}
\end{align*}

% -----------------------------------------------------------------------
% ROTATIONAL SYMMETRIES
% -----------------------------------------------------------------------
\subsection{Rotational Symmetries}

For an $n$th polygonal prism, there are $2n$ rotational symmetries

% -----------------------------------------------------------------------
% BASKETBALL PLAYERS
% -----------------------------------------------------------------------
\subsection{Basketball}

There are 18 different basketball players. How many ways can they organize themselves into 3
games each with 3 players versus 3 players?

$$\frac{18!}{(3!3!3!3!3!3!)(2!2!2!)3!}$$

${18 \choose 3} {15 \choose 3}$ double counts, so we need to divide by $2!$. We need to divide by $3!$ too because each game
is treated as distinct

% -----------------------------------------------------------------------
% INTEGER EQUATION BOUNDED
% -----------------------------------------------------------------------
\subsection{Integer Equation Bounded}

\begin{align*}
    &|U| = | \text{all vars } \leq 7 | + | \text{at least 1 var } \geq 8 | \\
    &| \text{all vars } \leq 7 | = |U| - | \text{at least 1 var } \geq 8 | \\
\end{align*}

$$|U| = {35 \choose 5}$$

$$|\text{at least 1 var being at least 8}| = |\text{1 var at least 8}| + |\text{2 vars at least 8}| + \ldots + |\text{6 vars at least 8}|$$

For each subset, "$p$ var(s) at least 8", we must multiply the arrangements by ${6 \choose p}$ since we need to choose where the "at least" variables are

For $p = 1$

\begin{align*}
    (a_1 + 8) + a_2 + a_3 + a_4 + a_5 + a_6 &= 30 \\
    a_1 + a_2 + a_3 + a_4 + a_5 + a_6 &= 22 \\
    {6 \choose 1}{27 \choose 5} &
\end{align*}

For $p = 2$

\begin{align*}
    (a_1 + 8) + (a_2 + 8) + a_3 + a_4 + a_5 + a_6 &= 30 \\
    a_1 + a_2 + a_3 + a_4 + a_5 + a_6 &= 14 \\
    {6 \choose 2}{19 \choose 5} &
\end{align*}

For $p = 3$

\begin{align*}
    (a_1 + 8) + (a_2 + 8) + (a_3 + 8) + a_4 + a_5 + a_6 &= 30 \\
    a_1 + a_2 + a_3 + a_4 + a_5 + a_6 &= 6 \\
    {6 \choose 3}{11 \choose 5} &
\end{align*}

And for $p = 4$ to $p = 6$, there's 0 solutions since the RHS would be negative

But we need to be careful because we would overcount if we just add up all the sets. The solutions for $p = 2$ are INCLUDED in $p = 1$, and same for
$p = 3$, so we need to apply Principle of Inclusion Exclusion. 

Recall that for PIE, if you're accounting for an ODD number of sets, then you subtract, but EVEN you add

so we add $p = 1$ and $p = 3$, but subtract $p = 2$

Doing all of this, we get

\begin{align*}
    &|\text{at least 1 var being at least 8}| = \\
    &{6 \choose 1}{27 \choose 5} - {6 \choose 2}{19 \choose 5} + {6 \choose 3}{11 \choose 5}
\end{align*}

And plugging this back into our base equation we get

$$| \text{all vars } \leq 7 | = \boxed{{35 \choose 5} - {6 \choose 1}{27 \choose 5} + {6 \choose 2}{19 \choose 5} - {6 \choose 3}{11 \choose 5}}$$

\end{multicols}

\end{document}